% ##############################################################################
\section{How an intelligence market creates value}
\label{sec:market_design}

To motivate the design choices made in the rest of the paper, we step back
and consider the participants, roles, and the central design difficulty of
a market for machine intelligence.

% ==============================================================================
\subsection{Market participants}
\label{sec:market_participants}

The \textbf{demand side} of the market is made of applications and autonomous agents
that consume LLM inference as an input: a customer-support agent, a
document-summarization pipeline, a coding assistant, or a multi-agent system
delegating a sub-task. These participants differ along several dimensions:

\begin{itemize}
  \item How much they value low latency.
  \item How much throughput they need.
  \item How much model intelligence (capability tier) their tasks require.
  \item How much they value a low error rate.
  \item How price-sensitive they are.
  \item How much of their traffic is predictable in advance versus bursty.
\end{itemize}

The \textbf{supply side} is made of model and compute providers: foundation-model
API vendors, open-weight model hosts, and resellers of spare inference
capacity. Sellers differ along several dimensions:

\begin{itemize}
  \item The capability tier of the models they can serve.
  \item Their typical latency, throughput, and reliability under load.
  \item Their marginal cost of serving a request, which itself depends on
    hardware utilization and contractual commitments elsewhere.
\end{itemize}

As in any market, the opportunity to trade arises from these differences:
a buyer with a latency-insensitive, cost-sensitive batch job and a seller
with idle capacity outside of peak hours can both be made better off by a
contract that a rigid, single posted price cannot express.

Figure~\ref{fig:two_sided_market} lays out both sides and the dimensions
along which their participants differ: demand-side participants differ
along latency and throughput needs, model-intelligence requirements, error
tolerance, price sensitivity, and traffic predictability, while supply-side
participants differ along capability tier, speed and reliability under load,
and marginal cost, and \NoesisMarket{} is the meeting point where bids and
asks expressing these differences are matched into contracts. Each side's
dimensions are drawn as short axes radiating from its box.

% rendered_images:begin
% ```tikz
% \definecolor{demandfill}{RGB}{246,225,232}
% \definecolor{demandline}{RGB}{107,42,68}
% \definecolor{supplyfill}{RGB}{251,235,212}
% \definecolor{supplyline}{RGB}{107,69,23}
% \definecolor{marketfill}{RGB}{211,227,243}
% \definecolor{marketline}{RGB}{31,78,121}
% \definecolor{axiscolor}{RGB}{163,177,192}
% \fontfamily{phv}\selectfont
% 
% \node[draw=demandline, rounded corners, fill=demandfill, minimum width=3.0cm,
%       minimum height=1.05cm, font=\normalsize\bfseries, align=center, text=demandline]
%   (demand) at (0,0) {Demand side\\[1pt]\small\normalfont Applications / agents};
% \node[draw=marketline, rounded corners, fill=marketfill, minimum width=2.9cm,
%       minimum height=1.05cm, font=\normalsize\bfseries, align=center, text=marketline]
%   (market) at (4.6,0) {NoesisMarket\\[1pt]\small\normalfont Meeting point};
% \node[draw=supplyline, rounded corners, fill=supplyfill, minimum width=3.3cm,
%       minimum height=1.05cm, font=\normalsize\bfseries, align=center, text=supplyline]
%   (supply) at (9.2,0) {Supply side\\[1pt]\small\normalfont Model \& compute providers};
% 
% \draw[->, >=stealth, thick, marketline] (demand.east) -- (market.west)
%   node[midway, above, font=\small] {Bids};
% \draw[->, >=stealth, thick, marketline] (supply.west) -- (market.east)
%   node[midway, above, font=\small] {Asks};
% 
% % Demand-side dimensions, fanned as short axes to the left of the box.
% \foreach \ang/\lab in {90/{traffic\\predictability}, 135/{latency and\\throughput},
%                        180/{capability\\required}, 225/{error\\tolerance},
%                        270/{price\\sensitivity}} {
%   \draw[->, >=stealth, thin, axiscolor] (demand.\ang) -- ++(\ang:1.5cm);
%   \path (demand.\ang) -- ++(\ang:2.05cm)
%     node[font=\scriptsize, align=center, text=demandline] {\lab};
% }
% 
% % Supply-side dimensions, fanned as short axes to the right of the box.
% \foreach \ang/\lab in {45/{capability\\tier}, 0/{speed and\\reliability},
%                        -45/{marginal\\cost}} {
%   \draw[->, >=stealth, thin, axiscolor] (supply.\ang) -- ++(\ang:1.5cm);
%   \path (supply.\ang) -- ++(\ang:2.05cm)
%     node[font=\scriptsize, align=center, text=supplyline] {\lab};
% }
% ```
% label=fig:two_sided_market
% caption=The two-sided \Noesis{} market: demand and supply meeting at \NoesisMarket{}, each side's dimensions drawn as short axes.
% rendered_images:end
% render_images:begin
\begin{figure}[H]
  \includegraphics[width=\linewidth]{02_market_design.tex.figs/02_market_design.1.png}
  \caption{The two-sided \Noesis{} market: demand and supply meeting at \NoesisMarket{}, each side's dimensions drawn as short axes.}
  \label{fig:two_sided_market}
\end{figure}
% render_images:end

% ==============================================================================
\subsection{Roles of an intelligence market}
\label{sec:market_roles}

An intelligence market, like any financial exchange, has to perform several
functions beyond simply connecting a buyer to a seller:

\begin{itemize}
  \item \textbf{Price discovery}: provide a common meeting ground where
    buyers and sellers can express demand and supply without prior
    bilateral negotiation.
  \item \textbf{Standardization}: establish a standardized contract schema
    (Section~\ref{sec:contracts_and_notation}) so that a bid and an ask
    submitted by unrelated parties are directly comparable.
  \item \textbf{Oversight}: meter whether a matched contract was actually
    fulfilled, a role played in \Noesis{} by \NoesisServer{} rather than
    by a separate compliance function.
  \item \textbf{Data generation}: generate reusable market data: clearing
    prices per class and round, and, as a byproduct of fulfillment
    monitoring, a record of measured provider performance over time.
  \item \textbf{Dissemination}: disseminate that clearing-price data in
    real time to participants and downstream systems, through the
    pricing-dissemination component introduced in
    Section~\ref{sec:modularity}, rather than requiring them to poll
    \NoesisMarket{} or wait for the next round.
\end{itemize}

As in traditional exchanges, we expect the market to be compensated for this
value creation through a small fee on cleared volume; we do not formalize a fee
schedule in this paper.

% ==============================================================================
\subsection{The bundling problem: quoting quality, not just quantity}
\label{sec:bundling}

The central design difficulty is choosing what, precisely, is being bought and
sold. A market for a homogeneous good only needs to clear on price and
quantity. Machine intelligence is not homogeneous: the same nominal price per
request can correspond to wildly different quality, and the same underlying
compute can be used to answer the same request at very different quality
depending on the model and routing policy used to serve it.

This is analogous to the base/quote asymmetry that a token-exchange protocol
must resolve before a limit order is even well-defined: a limit order only
makes sense once one has fixed which token plays the role of quantity and which
plays the role of price. \Noesis{} faces a similar prior question: before a bid
or an ask can be written down, the protocol must fix what a unit of
``intelligence'' is (Definition~\ref{def:task}) and what accompanies that
unit as a quality specification. We adopt the position that quality should
be expressed as two terms: a \emph{service class}, which says what each
request should get, and a \emph{reliability}, which says how often it gets
it. A service class has three largely orthogonal dimensions, each normalized
into a few discrete bands:

\begin{itemize}
  \item \textbf{Capability}: how good the answer is expected to be, a
    ``virtual intelligence'' level (e.g., low, medium, high). Bands are
    discrete rather than continuous scores, since bands are what providers
    can credibly commit to and what buyers can reason about.
  \item \textbf{Latency}: how long one request may take (e.g., fast, slow,
    or batch, for slow asynchronous jobs completed within 24 hours).
  \item \textbf{Throughput}: how fast one request produces output (e.g.,
    low, medium, high).
\end{itemize}

Reliability is then the fraction of requests, over the life of a contract,
that meet the contracted class in every dimension. It merges the three
per-request checks into a single measure. The protocol fixes only this
structure; the thresholds behind each band are set by a concrete
implementation.

Figure~\ref{fig:bundle_axes} makes this contrast visually concrete:
\Noesis{}'s traded unit is a bundle of a service class and a reliability,
contrasted with a conventional single-scalar good quoted purely on price per
token.

% rendered_images:begin
% ```tikz
% \definecolor{goodfill}{RGB}{223,237,224}
% \definecolor{goodline}{RGB}{143,183,154}
% \definecolor{goodtext}{RGB}{46,90,61}
% \definecolor{axiscolor}{RGB}{163,177,192}
% \definecolor{dividercolor}{RGB}{199,208,218}
% \fontfamily{phv}\selectfont
% \node[draw=goodline, rounded corners, fill=goodfill, minimum width=2.8cm, minimum height=1cm, font=\small, align=center, text=goodtext] (bundle) at (0,0) {traded unit\\(service class,\\reliability)};
% \draw[->, >=stealth, thick, axiscolor] (bundle.north) -- ++(0,1.1) node[above, font=\footnotesize, align=center] {Reliability};
% \draw[->, >=stealth, thick, axiscolor] (bundle.east) -- ++(1.2,0.7) node[right, font=\footnotesize, align=center] {Capability};
% \draw[->, >=stealth, thick, axiscolor] (bundle.east) -- ++(1.2,0) node[right, font=\footnotesize, align=center] {Latency};
% \draw[->, >=stealth, thick, axiscolor] (bundle.east) -- ++(1.2,-0.7) node[right, font=\footnotesize, align=center] {Throughput};
% \node[above, font=\footnotesize\bfseries] at (0,2.35) {service class + reliability};
% \draw[dashed, dividercolor] (5.3,-2.2) -- (5.3,2.4);
% \node[draw=goodline, rounded corners, fill=goodfill, minimum width=2.8cm, minimum height=1cm, font=\small, align=center, text=goodtext] (scalar) at (7.9,0) {traded unit\\(price per token)};
% \draw[->, >=stealth, thick, axiscolor] (scalar.east) -- ++(1.2,0) node[right, font=\footnotesize, align=center] {Price};
% \node[above, font=\footnotesize\bfseries] at (7.9,2.35) {single-scalar good};
% ```
% label=fig:bundle_axes
% caption=A service class and reliability bundle versus a single-scalar good.
% rendered_images:end
% render_images:begin
\begin{figure}[H]
  \includegraphics[width=\linewidth]{02_market_design.tex.figs/02_market_design.2.png}
  \caption{A service class and reliability bundle versus a single-scalar good.}
  \label{fig:bundle_axes}
\end{figure}
% render_images:end

Section~\ref{sec:contracts_and_notation} turns this into a formal contract. The
consequence of this design choice is that the market cannot clear on price
alone: two asks at the same price but different classes are not
equivalent, and matching must first establish compatibility on service class
and reliability before price discovery can proceed
(Section~\ref{sec:compatibility}).

% ==============================================================================
\subsection{Related commodity markets}
\label{sec:related_commodity_markets}

Electricity markets are the closest existing analogy. Electricity cannot be
usefully stored at scale, demand and supply must be balanced continuously, and
quality (voltage, frequency) is metered rather than taken on faith. Electricity
markets address this with a layered design: day-ahead markets clear the bulk of
expected demand a day in advance, real-time markets clear residual imbalances
close to delivery, and separate capacity markets pay generators for being
available rather than for energy delivered, penalizing non-delivery when it
matters most \cite{stoft2002, cramton2017}.

\Noesis{}'s per-class call auction (Section~\ref{sec:noesis_market}) borrows the
day-ahead/real-time separation as a candidate design (discussed further as an
open question in Section~\ref{sec:open_questions}), and its reputation and
pricing feedback loop (Section~\ref{sec:reputation}) borrows directly from
capacity-market non-performance penalties \cite{cramtonstoft2005}.

The analogy is imperfect in one important respect: electricity's quality
attributes (voltage, frequency) are physical and essentially free to meter,
while machine intelligence's capability attribute is not directly observable
and can only be estimated from the content of a response. This is the problem
that \NoesisServer{}'s fulfillment monitoring
(Section~\ref{sec:fulfillment_monitoring}) has to solve.
